\section{Distributional Calculus}
\label{app:calculus}
All statements use the finite component sets and measurable kernels of
Definition~\ref{def:sfd}. Exact atoms are Dirac laws; regularized blobs can
be represented among the continuous residual components. Component labels
are latent and need not partition the observed space.

\begin{proposition}[Normalization and latent construction]
\label{prop:normalization}
The SFD law is the marginal of the following construction: draw a component
$Z$ with probabilities $\pi_c$; given an atom, set $X=\xi_a$; given a branch,
draw $S$ with density $h_b$ and then $X\sim K_b(\cdot\mid S)$; given a
residual, draw $X$ with density $q_r$. In particular, $P_\theta(\mathcal X)=1$.
\end{proposition}
\begin{proof}
Each branch integral is a measure: for disjoint measurable sets,
monotone convergence passes their sum through the nonnegative kernel
integral. It has total mass $\int h_b(s)K_b(\mathcal X\mid s)\,ds=1$.
The finite nonnegative mixture therefore has total mass $\sum_c\pi_c=1$.
Conditioning the stated latent construction on $Z$ gives exactly its
probability on every measurable set, proving equality of laws.
\end{proof}

\begin{proposition}[Mixed-density evaluation]
\label{prop:mixed-density}
For fixed atom locations and fixed exact branch maps, a common reference is
\[
 \nu=\lambda_d+\sum_a\delta_{\xi_a}
            +\sum_{b\in\mathcal B_{\mathrm{ex}}}(\gamma_b)_\#ds,
\]
where $ds$ is restricted to the finite-volume coordinate domain. If $P_c$
denotes each normalized component law, then
\[
 p_\theta=\sum_c\pi_c\frac{dP_c}{d\nu}.
\]
For a continuous component with Lebesgue density $g_c$, its common-reference
density is $g_c\rho$, where $\rho=d\lambda_d/d\nu$. For an injective exact
branch, put $\nu_{\Gamma,b}=(\gamma_b)_\#ds$. Its contribution is
\[
 \pi_b h_b\{\gamma_b^{-1}(x)\}
                    \frac{d\nu_{\Gamma,b}}{d\nu}(x).
\]
Thus overlapping reference contributions must be converted before addition.
\end{proposition}
\begin{proof}
The branch reference measures are finite and their finite sum with Lebesgue
measure is sigma-finite. A coordinate null set has zero probability under
$h_b(s)ds$, so its pushforward probability is dominated by its pushforward
reference. All components are therefore dominated by $\nu$. Linearity of
the Radon--Nikodym derivative gives the mixture formula; its chain rule gives
the continuous and injective-branch expressions. In particular, a continuous
density contributes zero mass at an exact atom, even if its Lebesgue density
there is positive.
\end{proof}

For a regular immersion of intrinsic dimension $m_b$ with finite preimages,
the area formula gives, at almost every point of its image,
\[
 \frac{dP_b}{d\mathcal H^{m_b}}(x)
   =\sum_{s:\gamma_b(s)=x}\frac{h_b(s)}{J_b(s)},\qquad
 J_b(s)=\sqrt{\det\{D\gamma_b(s)^\top D\gamma_b(s)\}}>0.
\]
Relative to the branch's own pushforward reference, the density is the ratio
of this sum to $\sum_{s:\gamma_b(s)=x}1/J_b(s)$, not the unweighted sum of
coordinate densities. The common-reference conversion still applies.

\subsection{Expectation, Moments and Covariance}
For completeness, Proposition~\ref{prop:expectation} follows first for
nonnegative $g$ by conditioning on $Z$, then $S$, and using Tonelli's theorem.
For integrable signed $g$, apply the same argument to $g^+$ and $g^-$ and
subtract. This proves the full identity without a Lebesgue-only assumption.

In the Gaussian-tube case, whenever the second moments exist,
\begin{align*}
 \mu=\E_\theta X
 &=\sum_a\pi_a\xi_a+\sum_b\pi_b\int h_b(s)\gamma_b(s)\,ds
                         +\sum_r\pi_r\E_{q_r}X,\\
 \E_\theta(XX^\top)
 &=\sum_a\pi_a\xi_a\xi_a^\top
   +\sum_b\pi_b\int h_b(s)
       \{\gamma_b(s)\gamma_b(s)^\top+\Sigma_b(s)\}\,ds
   +\sum_r\pi_r\E_{q_r}(XX^\top).
\end{align*}
The covariance is $\E_\theta(XX^\top)-\mu\mu^\top$. A branch covariance
therefore separates variation of $\gamma_b(S)$ from the average conditional
noise covariance $\int h_b(s)\Sigma_b(s)\,ds$. If the component moment
generating functions and the following integrals are finite near the origin,
\[
 M_X(t)=\sum_a\pi_a e^{t^\top\xi_a}
 +\sum_b\pi_b\int h_b(s)
   e^{t^\top\gamma_b(s)+t^\top\Sigma_b(s)t/2}\,ds
 +\sum_r\pi_r M_{q_r}(t).
\]
Both formulas are choices of $g$ in the master identity.

\subsection{Marginal and Conditional Laws}
Let $J$ be a nonempty proper subset of coordinate indices, put $Y=X_{-J}$,
and write $T_J$ for projection onto $X_J$.

\begin{proposition}[Marginal law]
\label{prop:marginal}
For arbitrary SFD kernels,
\[
 (T_J)_\#P_\theta
  =\sum_a\pi_a\delta_{\xi_{a,J}}
    +\sum_b\pi_b\int h_b(s)(T_J)_\#K_b(\cdot\mid s)\,ds
    +\sum_r\pi_r(T_J)_\#Q_r.
\]
For Gaussian tubes, the projected branch kernel is
$\Normal(\gamma_{b,J}(s),\Sigma_{b,JJ}(s))$. Exact branches instead use
their projected pushforward laws, which may change dimension or acquire atoms.
\end{proposition}
\begin{proof}
Evaluate the original law on $T_J^{-1}(A)$ for a measurable set $A$.
Projection commutes with the finite sum and nonnegative kernel integral.
The Gaussian assertion is the marginal formula for a multivariate normal.
\end{proof}

\begin{proposition}[Gaussian branch conditional law]
\label{prop:gaussian-conditional}
For a full-rank Gaussian branch kernel, conditional on $(Z=b,S=s,Y=y)$,
$X_J$ is Gaussian with mean and covariance
\begin{align*}
 m_{J\mid-J,b}(s;y)
 &=\gamma_{b,J}(s)+\Sigma_{b,J,-J}(s)\Sigma_{b,-J,-J}(s)^{-1}
                         \{y-\gamma_{b,-J}(s)\},\\
 V_{J\mid-J,b}(s)
 &=\Sigma_{b,JJ}(s)-\Sigma_{b,J,-J}(s)
            \Sigma_{b,-J,-J}(s)^{-1}\Sigma_{b,-J,J}(s).
\end{align*}
\end{proposition}
\begin{proof}
Positive definiteness makes the conditioning covariance block invertible.
Completing the square in the joint Gaussian density gives the displayed
mean and Schur-complement covariance.
\end{proof}

\begin{proposition}[Full conditional mixture]
\label{prop:full-conditional}
Choose a common projected reference $\nu_{-J}$ that dominates the projected
component laws and, for each regularized branch, its projected slice laws.
Let $g_{a,-J}$, $g_{b,-J}(y\mid s)$ and $g_{r,-J}$ be their densities
relative to this same reference, and let $p_{\theta,-J}$ be the total density.
For $P_{\theta,-J}$-almost every $y$ with positive density, put
\[
 \tau_a(y)=\frac{\pi_a g_{a,-J}(y)}{p_{\theta,-J}(y)},\qquad
 \tau_b(s\mid y)=\frac{\pi_bh_b(s)g_{b,-J}(y\mid s)}{p_{\theta,-J}(y)},
 \qquad
 \tau_r(y)=\frac{\pi_r g_{r,-J}(y)}{p_{\theta,-J}(y)}.
\]
For regularized branches and exact atoms, the full conditional law is
\begin{align*}
 \Pbb(X_J\in A\mid Y=y)
 &=\sum_a\tau_a(y)\mathbf 1_A(\xi_{a,J})\\
 &\quad+\sum_b\int \tau_b(s\mid y)
       K_{b,J\mid-J}(A\mid s,y)\,ds
       +\sum_r\tau_r(y)Q_{r,J\mid-J}(A\mid y).
\end{align*}
For Gaussian tubes the branch conditional is given by
Proposition~\ref{prop:gaussian-conditional}. Exact branches whose slice laws
lack common domination use a posterior coordinate measure obtained by
disintegration, rather than the displayed density relative to $ds$.
\end{proposition}
\begin{proof}
Disintegrate the joint latent-observed law with respect to $Y$. The three
Bayes numerators are its atom, branch-coordinate and residual contributions;
their sum/integral is $p_{\theta,-J}$. They therefore normalize to one.
Integrating each component's conditional law against these posterior weights
gives the assertion. The exceptional set is projected-law null, as required
for regular conditional distributions.
\end{proof}

For Gaussian slices under a mixed projected reference,
$g_{b,-J}(y\mid s)=\phi_{d-|J|}\{y;\gamma_{b,-J}(s),
\Sigma_{b,-J,-J}(s)\}\rho_{-J}(y)$, with
$\rho_{-J}=d\lambda_{d-|J|}/d\nu_{-J}$. That conversion also applies to
continuous residuals. Conditional simulation uses $\tau_b(s\mid y)$,
whereas full-observation EM uses $\tau_b(s\mid x)$; they condition on
different observations.

\subsection{Special-Case Embeddings}
\label{app:special-cases}
\begin{proposition}[Special-case status]
\label{prop:special-cases}
The following restrictions give exact equalities of laws or estimators.
\begin{enumerate}
\item Atoms alone give finite discrete distributions. An atom and a specified
continuous residual give the corresponding atom-inflated mixture.
\item Residual-only SFDs contain normal, product, heavy-tailed and finite
Gaussian-mixture laws by choosing the corresponding normalized $q_r$.
\item A residual-only mixture of $n$ fixed data-centered Gaussian kernels
with weights $1/n$ is exactly that KDE estimator.
\item One Gaussian tube gives Proposition~\ref{prop:single-tube}; specified
multiple tubes give mixtures of those laws. A Gaussian factor-analyzer
distribution enters as its integrated residual Gaussian density.
\item Exact kernels $\delta_{\gamma_b(s)}$ give coordinate-pushforward
laws under an intrinsic or mixed reference, not ambient Lebesgue densities.
\item A specified copula density with its marginal densities can be chosen
as a residual law. This does not identify SFD itself as a copula construction.
\end{enumerate}
A learned metric or temporal/path-functional model is an adjacent connection
until its extra structure is explicitly included in a normalized kernel or
dynamic extension.
\end{proposition}
\begin{proof}
Set inactive weights to zero and substitute the stated component laws in
Definition~\ref{def:sfd}. Each retained law is normalized. The KDE assertion
holds conditional on its chosen centers and bandwidth; it is an estimator
equality, not an additional parametric generating model. The pushforward
case is interpreted using Proposition~\ref{prop:mixed-density}. Extra
metric or temporal structure is not implied by a static mixture restriction.
\end{proof}
